\section{Definitions}\label{sec:tsps-definitions}

A threshold structure-preserving signature scheme distributes the signing key of
an SPS scheme $\SPS$ between $n$ signers, so that any $t+1$ of them can jointly produce a signature, and no
$t$ of them can produce one alone.

Signers communicate over authenticated point-to-point channels, together with a
broadcast channel. A verifier holds the public key, and needs to know neither
$n$, nor $t$, nor which signers took part. Signers may be malicious, and may
collude to forge signatures or to recover the signing key. We write $\Cor$ for
the set of corrupted signers and $\mathcal{H}$ for its complement. We consider
static corruption: the adversary chooses which signers to corrupt before key
generation, subject to $|\Cor| \leq t$. \Cref{sec:tsps-adaptive} extends this to
adaptive corruption, where the adversary chooses them during protocol
execution.

\subsection{Structure-preserving signature schemes}\label{sec:tsps-sps}

Structure-preserving signature (SPS) schemes are signature schemes in which the messages,
public keys, and signatures are all group elements, and verification consists
solely of checking pairing product
equations~\cite{C:AFGHO10,C:AGHO11,TCC:AGOT14,C:KilPanWee15,AC:Groth15,RSA:Ghadafi16}.


An SPS scheme $\SPS$ over the bilinear groups of \Cref{sec:prelims-groups}
consists of the following
algorithms:

\begin{description}

  \item[$(\sk, \pk) \sample \SPSgen(\secparam, \ell)$ :] The key generation
algorithm, on input the security parameter and a vector length $\ell$, outputs
secret signing key $\sk$ and public verification key $\pk$, which consists of
group elements.

  \item[$\SPSsig \sample \SPSsign(\sk, \vec{\msg})$ :] The signing algorithm
signs a message $\vec{\msg}$ of $\ell$ group elements using secret key $\sk$, and
outputs
a signature $\SPSsig$ consisting of elements of $\group{1}$ and $\group{2}$.

  \item[$\{0,1\} \leftarrow \SPSverify(\pk, \vec{\msg}, \SPSsig)$ :]
The verification algorithm evaluates pairing product equations in $\pk$,
$\vec{\msg}$ and $\SPSsig$, and outputs $0$ or $1$.

\end{description}
%
Some schemes are rerandomisable: they also provide an algorithm
$\SPSrandomise$, which, without the secret key, turns a signature $\SPSsig$
into a new signature $\SPSsig'$ on the same message, distributed as a freshly
generated one.

A rerandomisable scheme cannot be strongly unforgeable, since rerandomising a
signature is itself a strong forgery.

Correctness and EUF-CMA security are defined as for digital signature schemes.
Messages may be elements of either source group. We consider schemes that sign
messages in $\group{1}^\ell$.

\subsection{Threshold structure-preserving signature schemes}

We write $\shr{x}$ for a sharing of $x$ between the $n$ signers, and $x_h$ for
the share of signer $h$. Signing is split
into an offline phase, run before the message is known, and an online phase. A
threshold SPS scheme $\TSPS$ consists of the following algorithms.

\begin{description}

  \item[{$(\shr{\sk}, \pk) \sample \TSPSgen(\secparam, \ell, t, n)$ :}] Run
jointly by the signers. On input the message length $\ell$, the corruption bound
$t$, and the number of signers $n > t$, gives each signer $h$ a share $\sk_h$
of the secret key, and outputs a public key $\pk$ together with the share
$\pk_h$ of it that each signer $h$ broadcasts.

  \item[{$\shr{\state} \sample \TSPSoffsign(\shr{\sk})$ :}] Run jointly by the
signers before the message is known. Gives each signer $h$ a share $\state_h$ of
the state for one future signature.

  \item[{$\SPSsig_h \leftarrow \TSPSonsign(\sk_h, \state_h, \vec{\msg})$ :}] Run
locally by signer $h$ on a message $\vec{\msg} \in \group{1}^\ell$. Outputs a
signature share $\SPSsig_h$.

  \item[{$\SPSsig \leftarrow \TSPScombine(\{\SPSsig_h\}_{h \in Q})$ :}] On input
signature shares from a set $Q$ of $t+1$ signers, outputs a signature
$\SPSsig$, without interaction.

  \item[$\{0,1\} \leftarrow \TSPSverify(\pk, \vec{\msg}, \SPSsig)$ :] Equal to
$\SPSverify(\pk, \vec{\msg}, \SPSsig)$.

\end{description}

This follows the syntax of threshold signatures with
preprocessing~\cite{C:BCKMTZ22}, except that our offline phase is a joint
protocol between the signers rather than a local step.

\paragraph{Correctness.} A threshold SPS scheme is correct if the following holds
for all $(\ell, t, n)$ with $n > t$, every $(\shr{\sk}, \pk)$ output by
$\TSPSgen(\secparam, \ell, t, n)$ and $\shr{\state}$ output by a run of
$\TSPSoffsign(\shr{\sk})$ that does not abort, every message $\vec{\msg} \in \group{1}^\ell$, and
every set $Q \subseteq [1,n]$ of $t+1$ signers:
%
\begin{equation*}
%
  \SPSverify\left(\pk, \vec{\msg}, \TSPScombine\left(\left\{
\TSPSonsign(\sk_h, \state_h, \vec{\msg}) \right\}_{h \in Q}\right)\right) = 1.
%
\end{equation*}

\paragraph{Unforgeability.} Threshold signatures have a hierarchy of
unforgeability notions, starting from TS-UF-0, which counts a forgery only if no
honest signer signed the message~\cite{C:BCKMTZ22}. TS-UF-1 counts a forgery
whenever fewer than $t+1-|\Cor|$ honest signers signed it, so that the adversary
never held $t+1$ shares on that message. The two notions are equivalent when the
adversary corrupts $t$ signers. TS-UF-1 is our target, as no signature should
exist unless enough honest signers agreed to that message. Both notions are
defined for non-interactive threshold signature schemes, with the adversary
collecting the signature shares~\cite{C:BCKMTZ22}.

\Cref{exp:ts-uf-1} gives the TS-UF-1 experiment of~\cite{C:BCKMTZ22} in the
notation of this chapter, with $t+1$ signers needed to sign rather than $t$.

\begin{experimentbox}{Threshold signature schemes: TS-UF-1~\cite{C:BCKMTZ22}.}{ts-uf-1}

  \begin{enumerate}

    \item The adversary outputs $\Cor \subseteq [1, n]$ with $|\Cor| \leq t$,
and the challenger sets $\mathcal{H} := [1, n] \setminus \Cor$. It runs $(\pk,
\algo{aux}, \sk_1, \ldots, \sk_n) \sample \TSgen^{\hash}(\secparam, t, n)$, stores
$\sk_h$, $\pk$ and $\algo{aux}$ in the state $\algo{st}_h$ of every $h \in
\mathcal{H}$, and sends $\pk$, $\algo{aux}$ and $\{\sk_h\}_{h \in \Cor}$ to
$\adv$.

    \item The adversary is given oracle access to the following:

      \begin{description}

	\item[$\mathcal{O}_{\algo{Pre}}(h)$:] If $h \notin \mathcal{H}$, $\adv$
loses. Otherwise run $(\algo{pp}, \algo{st}_h) \sample
\TSpre^{\hash}(\algo{st}_h)$ and return $\algo{pp}$.

	\item[$\mathcal{O}_{\algo{PSign}}(h, \vec{\msg}, Q)$:] If $Q \not\subseteq
[1, n]$ or $h \notin \mathcal{H} \cap Q$, $\adv$ loses. Otherwise run
$(\SPSsig_h, \algo{st}_h) \sample \TSpsign^{\hash}(\vec{\msg}, Q, h,
\algo{st}_h)$, add $h$ to $S(\vec{\msg})$ if $\SPSsig_h \neq \perp$, and return
$\SPSsig_h$.

	\item[$\mathcal{O}_{\hash}(x)$:] Return $\hash(x)$.

      \end{description}

    \item The adversary outputs $(\vec{\msg}^\star, \SPSsig^\star)$.

  \end{enumerate}

  $\adv$ wins the experiment if $\TSverify^{\hash}(\pk, \vec{\msg}^\star,
\SPSsig^\star) = 1$ and $|S(\vec{\msg}^\star)| < t + 1 - |\Cor|$.

\end{experimentbox}

\Cref{exp:tsps-ts-uf-1} adapts this experiment to our setting. Our offline
oracle runs $\TSPSoffsign$ between all signers, with $\adv$ acting for $\Cor$,
and returns the view of the corrupted signers. Our online oracle runs
$\TSPSonsign$ for one honest signer. There is no leader, since $\adv$ calls the
honest signers directly. There is no random oracle, since our schemes do not use
one. Each session is used for one message only, since reusing a state reuses the
signing randomness. The challenger runs $\TSPSgen$, idealising key generation as
in the existing definitions~\cite{C:BCKMTZ22,C:CriKomMal23}.

The advantage of a PPT adversary $\adv$ is defined as follows:
%
\begin{equation*}
%
  \Adv_{\TSPS, \adv}^{\mathsf{TS\text{-}UF\text{-}1}}(\secparam) = \Pr[\adv\text{ wins}],
%
\end{equation*}
%
with the experiment in \Cref{exp:tsps-ts-uf-1}. A threshold SPS scheme is TS-UF-1
secure if this advantage is negligible for every PPT adversary $\adv$.

\begin{experimentbox}{Threshold structure-preserving signature schemes: TS-UF-1}{tsps-ts-uf-1}

  \begin{enumerate}

    \item The adversary outputs $(n, t, \Cor)$ with $|\Cor| \leq t$, and the
challenger sets $\mathcal{H} := [1,n] \setminus \Cor$.

    \item The challenger generates $(\shr{\sk}, \pk) \sample \TSPSgen(\secparam,
\ell, t, n)$, and sends $\pk$, the shares $\{\pk_h\}_{h \in [1,n]}$ that the
signers broadcast to open it, and $\{\sk_h\}_{h \in \Cor}$ to $\adv$.

    \item The adversary is given oracle access to the following:

      \begin{description}

	\item[$\mathcal{O}_{\algo{Off}}(\mathsf{sid})$:] If $\mathsf{sid}$ has
been used before, return $\perp$. Otherwise run $\TSPSoffsign$ between all $n$
signers, with $\adv$ acting for $\Cor$, store the shares
$\shr{\state}_{\mathsf{sid}}$, and return the view of the corrupted signers. If
the run aborts, return $\perp$ and mark $\mathsf{sid}$ as used.

	\item[$\mathcal{O}_{\algo{On}}(h, \mathsf{sid}, \vec{\msg})$:] If $h
\notin \mathcal{H}$, $\mathsf{sid}$ was not initialised by
$\mathcal{O}_{\algo{Off}}$, or $\mathsf{sid}$ has been used with a message other
than $\vec{\msg}$, return $\perp$. Otherwise compute $\SPSsig_h \leftarrow
\TSPSonsign(\sk_h, \state_{\mathsf{sid},h}, \vec{\msg})$, add $h$ to
$S(\vec{\msg})$, and return $\SPSsig_h$.

      \end{description}

    \item The adversary outputs $(\vec{\msg}^\star, \SPSsig^\star)$.

  \end{enumerate}

  $\adv$ wins the experiment if $\SPSverify(\pk, \vec{\msg}^\star, \SPSsig^\star) = 1$
and $|S(\vec{\msg}^\star) \cap \mathcal{H}| < t + 1 - |\Cor|$, where
$S(\vec{\msg})$ records which signers have issued a signature share on
$\vec{\msg}$.

\end{experimentbox}

We prove that our construction is TS-UF-1 secure in
\Cref{sec:tsps-security}.

\subsection{Secure multiparty computation}\label{sec:tsps-mpc}

We describe our schemes as sequences of calls to interfaces of a secure
multiparty computation protocol. For a secret value $\alpha \in \Zp$ shared
between $n$ parties, such that any $t+1$ of them can reconstruct it by a
pre-defined linear combination, we write $\shr{\alpha}$ for the vector
$(\alpha_1, \ldots, \alpha_n)$ in which $\alpha_h$ is the share held by party
$h$. We use the same notation $\shr{a}$ for a sharing of a group element $a \in
\group{}$, where $\group{}$ is either source group.

A group element $a$ is shared in the group: the shares of $a$ are
themselves elements of $\group{}$, and reconstruction combines them with the
group operation, as a product of shares raised to fixed coefficients, so that
the linear operations below remain local.

We assume this notation incorporates the authentication components that
malicious security requires, which we otherwise abstract away. They mean that a
modified share is detected at opening. In protocol boxes, $\shr{\alpha}$ also
stands for the party's own share.

The protocol provides the following interfaces, each run by all parties
together. The interfaces over $\group{}$ exist once for each source group, and
we write $\algo{ExpGrp}_1$ for the one over $\group{1}$.

  \begin{description}

    \item[$\shr{\alpha} \sample \algo{RandShare}()$ :] Outputs a sharing of a
uniform $\alpha \sample \Zp$ that no party learns.

    \item[$\shr{\gamma} \sample \algo{Mul}(\shr{\alpha}, \shr{\beta})$ :] Outputs
a sharing of $\gamma = \alpha\beta$.

    \item[$\shr{\alpha^{-1}} \sample \algo{Inv}(\shr{\alpha})$ :] On input a sharing
of $\alpha \neq 0$, outputs a sharing of $\alpha^{-1}$.

    \item[$\shr{c} \sample \algo{ExpGrp}(\shr{\alpha}, \shr{b})$ :] On input a
sharing of $\alpha \in \Zp$ and a sharing of $b \in \group{}$, outputs a sharing
of $c = b^\alpha$.

    \item[$\alpha \leftarrow \algo{Open}(\shr{\alpha})$ :] Reconstructs $\alpha$
and gives it to every party, or to a designated one. Applies equally to a
sharing of a group element.

  \end{description}


Linear operations need no interaction, and are not exposed as interfaces. Given
sharings $\shr{\alpha}$, $\shr{\beta}$ and public $\phi, \psi, \theta \in \Zp$,
each party applies the combination to its own shares to obtain the following:
%
\begin{equation*}
%
  \shr{\phi\alpha + \psi\beta + \theta} = \phi\shr{\alpha} + \psi\shr{\beta} +
\theta.
%
\end{equation*}
%
In a group, given sharings $\shr{a}$, $\shr{b}$ and a public group element $d$,
each party obtains the following:
%
\begin{equation*}
%
  \shr{a^{\phi}b^{\psi}d} = \shr{a}^{\phi}\shr{b}^{\psi}d.
%
\end{equation*}
%
Under additive sharing, only one designated party adds the public $\theta$, or
multiplies by the public $d$, and the others leave their shares unchanged.
Raising a public group element to a shared exponent gives a
sharing of the result~\cite{LC:ADEO21}. For a public $d$, this is
the following:
%
\begin{equation*}
%
  \shr{d^{\alpha}} = d^{\shr{\alpha}}.
%
\end{equation*}
%

$\algo{ExpGrp}$ raises a shared group element to a shared exponent.

$\algo{Inv}$ is built from $\algo{Mul}$ and
$\algo{Open}$. The parties sample $\shr{\beta} \sample \algo{RandShare}()$, open
$\gamma = \alpha\beta$, which is uniform and so reveals nothing, and take
$\shr{\alpha^{-1}} = \gamma^{-1}\shr{\beta}$. Where a scheme samples from
$\Zp^*$, the parties sample from $\Zp$ and abort if an opened $\gamma$ is zero,
which happens with probability $1/p$. Where no inverse is taken, sampling from
$\Zp$ instead of $\Zp^*$ changes the distribution only negligibly.

Finally, we require the MPC protocol to be secure with abort against up to $t$
statically corrupted parties, in the standard simulation-based
sense~\cite{EPRINT:Lindell16}. Security with abort means that corrupted parties
may make a call return $\bot$ to every honest party, but cannot make it return
a wrong output. For every adversary controlling the corrupted parties, there is
a PPT simulator $\simu$ that interacts with the adversary in place of the honest
parties, given only the outputs of the corrupted parties and the values opened.
We write $\Adv^{\mathsf{mpc}}(\secparam)$ for the largest advantage of a PPT
distinguisher between a real execution of a sequence of interface calls and one
simulated by $\simu$.
